\documentclass[runningheads]{llncs}
\usepackage[T1]{fontenc}
\usepackage{lmodern}
\usepackage{microtype}
\usepackage{amsmath,amssymb}
\usepackage{booktabs,tabularx,array}
\usepackage{xurl}
\usepackage[hidelinks]{hyperref}
\emergencystretch=2em

\begin{document}
\title{GovMut-Health: A Mutation Benchmark for Sequence-Sensitive Governance Faults in Stateful Health-AI Systems}
\titlerunning{GovMut-Health}
\author{Nguyen Ngoc Thien \and Trinh Minh Quang}
\authorrunning{N. N. Thien and T. M. Quang}
\institute{Hai Ba Trung High School, Hue, Vietnam\\\email{kakangocthien109@gmail.com}}
\maketitle
\thispagestyle{empty}\pagestyle{empty}

\begin{abstract}
Stateful health-AI governance failures often depend on an ordered sequence of operations rather than an isolated malformed input. A system may behave correctly when a disclosure snapshot is compiled, consent is active, and a proposal is generated, yet still admit an unauthorized persistent write if patient consent is revoked, policy epochs advance, a stale proposal is retried, or un-invalidated cache rows are re-inflated before durability. While general mutation testing evaluates syntactic AST mutations (e.g., SWE-Mutation) and formal agent verification tests idealized state machines (e.g., MemTX), neither captures sequence-sensitive governance failures in stateful health architectures. GovMut-Health establishes the first domain-specific mutation benchmark for healthcare AI, mapping 15 multi-action defect classes directly to clinical data lifecycle invariants and cryptographic Merkle state transition verification. We formalize Theorem~3 (Cryptographic Merkle Digest State Transition Verification and Bound-Commit Security), proving that under standard collision resistance, proposals bound to stale or drifted states are rejected with overwhelming probability. In a sealed experiment evaluating 45 non-equivalent mutants across 16 method/seed slots (720 executions, zero infrastructure errors), mutation scores were 35.6\% for regression tests, 8.9\% for stateless properties, 13.3\% for state-machine testing, and 44.4\% for the combined strategy. The combined strategy detected all 16 mutants killed by regression and four additional sequence-sensitive mutants, demonstrating empirical complementarity on this frozen corpus.
\end{abstract}

\section{Introduction}
A stateful governance contract in healthcare AI can fail even when every individual API request satisfies local syntactic and schema preconditions. The defect typically requires an interleaved execution sequence: disclose patient context, modify consent or governing policy, retain a stale proposal, retry a previously conflicted write, and attempt durability. Standard unit and example-based regression suites frequently cover happy paths and known historical bugs, but they fail to systematically explore the combinatorial transition orderings where sequence-dependent governance defects emerge.

General mutation testing evaluates test suite quality by injecting syntactic faults into source code (e.g., SWE-Mutation~\cite{swemutation}), while transactional agent memory benchmarks test idealized key-value state machines (e.g., MemTX~\cite{memtx}). However, neither paradigm models the domain-specific invariants of longitudinal health data: bitemporal valid-time and transaction-time intervals, clinical terminology ontologies (LOINC, SNOMED CT, RxNorm), dynamic patient consent revocations, and cryptographic Merkle witness bindings.

GovMut-Health addresses this fundamental gap. It introduces an executable domain-specific mutation model for stateful health AI, defines 15 multi-action defect classes tied to clinical lifecycle invariants, formalizes cryptographic Merkle digest state transition verification (Theorem~3), and benchmarks four frozen testing strategies under a sealed experimental protocol.

\section{Background and Novelty Boundary}
\subsection{Model-Based, Property-Based, and Access-Control Testing}
QuickCheck introduced property-based random testing, and model-based testing has long utilized stateful behavioral specifications to generate and shrink execution traces~\cite{quickcheck,utting}. Mutation testing measures test adequacy by seeding controlled faults and evaluating kill ratios~\cite{jia}. In authorization systems, model-based generation and policy mutation have been applied to access-control matrices~\cite{pretschner2008,xu2012,nist800192}. GovMut-Health does not claim novelty for these foundational testing methods.

\subsection{Differentiation from SWE-Mutation and MemTX}
Recent software engineering benchmarks evaluate LLM-generated test suites on syntactic code mutations (SWE-Mutation~\cite{swemutation}), while stateful agent memory frameworks check transactional belief commits and snapshot isolation on abstract memory stores (MemTX~\cite{memtx}, MemTxn~\cite{memtxn}). Furthermore, bitemporal memory formalisms and verified concurrent agent runtimes have highlighted sequence-sensitive execution anomalies~\cite{toki,verifiedconcurrency}.

However, critical differences separate GovMut-Health from these prior approaches:
\begin{enumerate}
    \item \textbf{Syntactic AST vs.\ Domain-Specific Governance Mutations:} SWE-Mutation applies generic mutation operators (e.g., arithmetic operator replacement, conditional statement inversion, statement deletion) at the Abstract Syntax Tree (AST) level. These syntactic perturbations cannot synthesize multi-step domain errors such as dropping bitemporal conflict markers, bypassing dynamic consent revocations, or omitting cryptographic Merkle root recomputations. GovMut-Health specifically targets domain-level enforcement gates.
    \item \textbf{Idealized Key-Value Memory vs.\ Clinical Lifecycle Invariants:} MemTX and MemTxn model idealized transactional memory blocks where state updates are evaluated as isolated read-write registers. In contrast, health AI governance involves complex clinical ontologies, two-dimensional temporal cutoffs ($t_{\text{valid}} \le \text{valid\_cutoff} \land t_{\text{known}} \le \text{known\_cutoff}$), patient consent scopes, and immutable audit logs.
\end{enumerate}

\subsection{Surviving Contribution}
GovMut-Health contributes: (1) a formal domain-specific mutation taxonomy comprising 15 multi-action defect classes; (2) a cryptographic verification model for Governed State Transitions (Theorem~3); (3) a curated, dual-model reviewed corpus of 45 non-equivalent mutants mapped to contract invariants; and (4) an empirical mutant-level evaluation of four testing strategies on real healthcare API and persistence boundaries.

\section{Reference State Model and Cryptographic Verification}
\label{sec:state_model}
\subsection{Abstract State Formulation}
The reference governance state is formalized as a tuple:
\begin{equation}
X = (u, v_s, v_p, v_c, c, a, p, H, P, L, A),
\end{equation}
where $u \in \mathcal{U}$ is the patient subject; $v_s, v_p, v_c \in \mathbb{N}$ denote state, policy, and consent version coordinates; $c \in \{\text{OPT\_IN}, \text{OPT\_OUT}, \text{RESTRICTED}\}$ is the active consent status; $a \in \mathcal{A}$ is the acting principal; $p \in \mathcal{P}$ is the declared clinical purpose; $H$ is the registry of issued task-bounded snapshots; $P$ is the set of pending persistent proposals; $L$ represents the canonical ledger state; and $A$ is the append-only audit trail.

\begin{table}[t]
\caption{Core governance invariants mapped to lifecycle enforcement rules.}
\label{tab:inv}
\centering\scriptsize
\begin{tabularx}{\textwidth}{p{0.08\textwidth}X}
\toprule
ID & Invariant Specification \\
\midrule
P1 & \textbf{Idempotency:} Idempotent replay cannot duplicate committed persistent state.\\
P2 & \textbf{Causal Precedence:} A stale base version cannot overwrite a newer partition state ($V(e_k) == v_s(e_k)$).\\
P3 & \textbf{Policy Freshness:} A policy epoch advance invalidates any proposal generated under stale policy.\\
P4 & \textbf{Consent Freshness:} Consent revocation immediately invalidates dependent disclosure and write rights.\\
P5--P7 & \textbf{Identity Binding:} Subject, actor, and purpose coordinates cannot silently migrate or cross-bind.\\
P8 & \textbf{Snapshot Integrity:} Expired ($t > t_{\text{exp}}$) or Merkle digest-mismatched snapshots cannot authorize commit.\\
P9 & \textbf{Supersession:} A superseded assertion is strictly excluded from subsequent active disclosures.\\
P10 & \textbf{Bitemporal Coherence:} Historical reconstruction preserves bitemporal valid/knowledge truth and conflicts.\\
P11 & \textbf{Provenance Closure:} Every committed clinical assertion maintains an unbroken causal PROV-O trace.\\
P12 & \textbf{Ledger Primacy:} Canonical ledger state survives deletion or rebuild of derived caches/snapshots.\\
P13 & \textbf{Atomic Durability:} Failed commits abort atomically: zero partial state changes or corrupt audit rows.\\
P14 & \textbf{Audit Reconstructability:} Successful commits are fully reconstructable from the immutable audit trail.\\
P15 & \textbf{Authorized Retry:} Retries cannot turn a stale or unauthorized proposal into an admitted write.\\
\bottomrule
\end{tabularx}
\end{table}

\subsection{Cryptographic Merkle Digest State Transition Verification}
To prevent Time-of-Check to Time-of-Use (TOCTOU) exploits and context spoofing, disclosures are bound cryptographically. A Task-Bounded Health State Snapshot (THSS) $\mathcal{H}$ is compiled at disclosure timestamp $t_1 = t_{\text{disclose}}(\mathcal{H})$ over disclosed clinical evidence $E_H$, base entity partition version vector $\vec{v}_0$, policy epoch $\Phi_{\text{policy}}$, and consent coordinate $\Sigma_{\text{consent}}$:
\begin{equation}
\mathcal{H} = \operatorname{THSS}(u, a, r, \phi, \tau, \vec{v}_0, id_H, d_H, E_H),
\end{equation}
where the canonical cryptographic Merkle digest $d_H$ is computed as:
\begin{equation}
d_H = \operatorname{MerkleRoot}\left(E_H \cup \vec{v}_0 \cup \Phi_{\text{policy}} \cup \Sigma_{\text{consent}}\right).
\end{equation}

When an agent proposes a state update $P = \langle u, a, r, \phi, \tau, v_s, v_{\pi}, v_c, id_H, d_H, \Delta, \rho \rangle$ at commit timestamp $t_2 = t_{\text{commit}}(P) > t_1$, the Governed State Transition (GST) revalidates state freshness, governance epochs, and cryptographic Merkle integrity under row-level database locks.

\begin{theorem}[Cryptographic Merkle State Transition Non-Forgeability and Bound-Commit Security]
\label{thm:crypto_merkle}
Let $\mathcal{H}_{\text{hash}}$ be a cryptographic hash function modeled as a random oracle / collision-resistant hash family under EUF-CMA / IND-CCA security. Let $\mathcal{H}$ be an authorized THSS snapshot issued at $t_1$ with Merkle root digest $d_H$. For any probabilistic polynomial-time (PPT) adversary $\mathcal{A}$ or drifting model generating proposal $P$ at $t_2 > t_1$:
\begin{equation}
\Pr\left[\operatorname{GST\_Commit}(P, t_2) = \text{True} \;\middle|\; \exists e_k \in \operatorname{Deps}(P) : V(e_k)_{t_2} > v_s(e_k) \;\lor\; \Sigma_{t_2} \neq \Sigma_{t_1} \;\lor\; d_H \neq \operatorname{RecomputeMerkleRoot}(id_H)\right] \le \operatorname{negl}(\lambda),
\end{equation}
where $\lambda$ is the security parameter and $\operatorname{negl}(\lambda)$ is a negligible function.
\end{theorem}

\begin{proof}
Suppose for contradiction that $\operatorname{GST\_Commit}(P, t_2) = \text{True}$ while at least one security premise is violated:
\begin{enumerate}
    \item \textbf{Payload / Evidence Tampering:} Suppose the payload $E_H$ or coordinates $\vec{v}_0$ were mutated to $E'_H \neq E_H$. For the proposal to satisfy $\operatorname{Bound}(P, \mathcal{H})$, the adversary must find $E'_H$ such that $\operatorname{MerkleRoot}(E'_H \cup \dots) = d_H$. By the collision resistance of $\mathcal{H}_{\text{hash}}$, the probability of discovering a second pre-image or collision is bounded by $\operatorname{negl}(\lambda)$.
    \item \textbf{State Partition Drift:} If an intermediate transaction $\mathcal{T}_{\text{state}}$ committed at $t \in (t_1, t_2)$ modifying partition $e_k \in \operatorname{Deps}(P)$, the monotonic version counter was incremented: $V(e_k)_{t_2} = V(e_k)_{t_1} + 1 > v_s(e_k)$. Under atomic execution, GST acquires a row-level lock (\texttt{SELECT ... FOR UPDATE}) on $e_k$ and evaluates predicate $\operatorname{StateCurrent}(P)$. Because $V(e_k)_{t_2} > v_s(e_k)$, $\operatorname{StateCurrent}(P)$ strictly evaluates to $\text{False}$.
    \item \textbf{Governance / Consent Invalidation:} If consent was revoked ($\Sigma_{t_2} \neq \Sigma_{t_1}$) or policy epoch advanced, re-reading the committed governance state under the lock forces $\operatorname{GovernanceCurrent}(P)$ to evaluate to $\text{False}$.
\end{enumerate}
Because admission requires the logical conjunction of $\operatorname{Bound}(P, \mathcal{H}) \land \operatorname{StateCurrent}(P) \land \operatorname{GovernanceCurrent}(P)$, the proposal is unconditionally aborted. Thus, the probability of unauthorized commit is bounded by $\operatorname{negl}(\lambda)$. $\blacksquare$
\end{proof}

\section{Testing Strategies}
We evaluate four distinct testing strategies against the same system revision:
\begin{itemize}
    \item \textbf{M0 (Regression):} Hand-crafted unit, integration, and contract tests designed by domain developers to guard known regressions and standard API behaviors.
    \item \textbf{M1 (Stateless Property Testing):} Generates randomized inputs and parameter variations over isolated API endpoints, checking local invariants without preserving multi-step history.
    \item \textbf{M2 (State-Machine Testing):} An executable rule-based state machine that generates, explores, and shrinks multi-step transition histories biased toward precondition boundaries.
    \item \textbf{M3 (Combined Strategy):} The union of M0, M1, and M2 executed under a frozen overall compute budget.
\end{itemize}

\section{Controlled Mutation Benchmark: 15 Multi-Action Defect Classes}
\label{sec:benchmark}

\begin{table}[t]
\caption{Cross-Paradigm Mutation and Verification Benchmark Comparison.}
\label{tab:paradigm_comparison}
\centering\scriptsize
\begin{tabularx}{\textwidth}{l c c c c c}
\toprule
\textbf{Paradigm} & \textbf{State Invariants} & \textbf{Lineage Trace} & \textbf{Bitemporal Math} & \textbf{Merkle Root} & \textbf{Target Domain}\\
\midrule
Syntactic Mutants (MutPy / SWE-Mutation)~\cite{jia,swemutation} & $\times$ & $\times$ & $\times$ & $\times$ & Generic Source Code\\
MemTxn / Cordon~\cite{memtxn,cordon} & \checkmark & $\times$ & $\times$ & $\times$ & Agent Memory Stores\\
TOKI Algebra~\cite{toki} & \checkmark & $\times$ & \checkmark & $\times$ & Bitemporal Memory\\
FHIR Version OCC~\cite{fhirhttp} & \checkmark & $\times$ & $\times$ & $\times$ & EHR REST Endpoints\\
\textbf{GovMut-Health (Ours)} & \textbf{\checkmark} & \textbf{\checkmark} & \textbf{\checkmark} & \textbf{\checkmark} & \textbf{Stateful Health AI}\\
\bottomrule
\end{tabularx}
\end{table}

\subsection{15 Multi-Action Defect Classes}
Unlike traditional mutation testing where operators mutate single AST tokens, GovMut-Health defines **15 multi-action defect classes** ($C_1$--$C_{15}$). Each class represents a realistic governance defect whose observable violation manifests only across an ordered sequence of state-machine actions (e.g., Disclose $\to$ Mutate Governance $\to$ Generate Proposal $\to$ Commit/Retry):

\begin{enumerate}
    \item \textbf{$C_1$ (MUT1 -- Stale Base Version Overwrite):} Disables optimistic version check ($V(e_k) == v_s(e_k)$), allowing writes based on stale snapshots to overwrite concurrent updates (violating P2).
    \item \textbf{$C_2$ (MUT2 -- Policy Epoch Invalidation Bypass):} Omits policy revalidation at commit time, admitting writes generated under outdated clinical governance policies (violating P3, P15).
    \item \textbf{$C_3$ (MUT3 -- Dynamic Consent Revocation Bypass):} Omits commit-time consent checks, allowing writes after a patient has explicitly revoked access (violating P4, P15).
    \item \textbf{$C_4$ (MUT4 -- Subject Boundary Cross-Leak):} Bypasses patient identity equality ($u_P == u_H$), allowing proposals to persist across different patient records (violating P5).
    \item \textbf{$C_5$ (MUT5 -- Actor Coordinate Escalation):} Omits principal identity verification ($a_P == a_H$), admitting unauthenticated or unauthorized agent writes (violating P6).
    \item \textbf{$C_6$ (MUT6 -- Purpose Specification Creep):} Bypasses clinical purpose matching ($p_P == p_H$), permitting treatment-scoped data to be modified under research/commercial scopes (violating P7).
    \item \textbf{$C_7$ (MUT7 -- Expired Snapshot Leniency):} Bypasses snapshot TTL validation ($t > t_{\text{exp}}$), admitting proposals generated from arbitrarily expired disclosures (violating P8).
    \item \textbf{$C_8$ (MUT8 -- Cryptographic Merkle Digest Bypass):} Skips Merkle root recomputation and verification ($d_H(P) == \operatorname{RecomputeMerkleRoot}(id_H)$), permitting payload tampering (violating P8, Theorem~\ref{thm:crypto_merkle}).
    \item \textbf{$C_9$ (MUT9 -- Snapshot-to-Unbound Downgrade):} Permits a THSS-consuming proposal to strip its snapshot binding metadata and commit as an unconstrained write (violating P3--P8).
    \item \textbf{$C_{10}$ (MUT10 -- Transactional Audit/State Atomicity Fracture):} Commits persistent clinical state while dropping or corrupting the corresponding audit trail entry (violating P13, P14).
    \item \textbf{$C_{11}$ (MUT11 -- Provenance Graph Lineage Severance):} Drops causal PROV-O parent edges linking newly committed assertions to their source disclosure snapshot (violating P11, P14).
    \item \textbf{$C_{12}$ (MUT12 -- Stale Derived Snapshot / Cache Re-inflation):} Serves stale cached clinical data or rebuilds derived views from un-invalidated caches after underlying entity updates (violating P4, P9, P12).
    \item \textbf{$C_{13}$ (MUT13 -- Superseded Assertion Re-selection):} Selects historically superseded clinical records as active medical truths during snapshot compilation (violating P9, P10).
    \item \textbf{$C_{14}$ (MUT14 -- Idempotency Key Replay Collision):} Disables idempotency key uniqueness constraints, allowing duplicate request executions to produce duplicate clinical side-effects (violating P1).
    \item \textbf{$C_{15}$ (MUT15 -- Unauthorized Conflict Retry Escalation):} Allows a proposal rejected due to state conflict to retry and succeed without re-reading state or refreshing governance authorization (violating P15).
\end{enumerate}

\begin{table}[t]
\caption{Prespecified 15 multi-action defect classes and invariant mappings.}
\label{tab:mut}
\centering\scriptsize
\begin{tabularx}{\textwidth}{p{0.07\textwidth}p{0.28\textwidth}p{0.36\textwidth}X}
\toprule
Class & Defect Name & Multi-Action Fault Sequence & Invariant \\
\midrule
MUT1 & Stale Base Overwrite & Disclose($v_1$) $\to$ Commit($v_2$) $\to$ CommitProposal($v_1$) & P2\\
MUT2 & Policy Drift Bypass & Disclose($\Phi_1$) $\to$ AdvancePolicy($\Phi_2$) $\to$ CommitProposal($\Phi_1$) & P3, P15\\
MUT3 & Consent Revoke Bypass & Disclose($\Sigma_1$) $\to$ RevokeConsent($\Sigma_2$) $\to$ CommitProposal($\Sigma_1$) & P4, P15\\
MUT4 & Subject Cross-Leak & Disclose($u_A$) $\to$ Propose($u_A$) $\to$ CommitProposal($u_B$) & P5\\
MUT5 & Actor Escalation & Disclose($a_A$) $\to$ Propose($a_A$) $\to$ CommitProposal($a_B$) & P6\\
MUT6 & Purpose Creep & Disclose($p_1$) $\to$ Propose($p_1$) $\to$ CommitProposal($p_2$) & P7\\
MUT7 & Expired Snapshot & Disclose($t_{\text{exp}}$) $\to$ AdvanceClock($>t_{\text{exp}}$) $\to$ CommitProposal & P8\\
MUT8 & Merkle Digest Tamper & Disclose($d_H$) $\to$ MutatePayload($E'_H$) $\to$ CommitProposal($d_H$) & P8, Thm~3\\
MUT9 & Unbound Downgrade & Disclose($\mathcal{H}$) $\to$ StripBinding($\mathcal{H}$) $\to$ DirectCommit & P3--P8\\
MUT10 & Atomicity Fracture & CommitProposal $\to$ PersistState $\to$ DropAuditEntry & P13, P14\\
MUT11 & Lineage Severance & CommitProposal $\to$ DropProvEdge($id_H$) $\to$ PersistState & P11, P14\\
MUT12 & Cache Re-inflation & CommitUpdate $\to$ InvalidateCache[Skip] $\to$ ReadStaleDerived & P4, P9, P12\\
MUT13 & Superseded Re-selection & Supersede($A_1 \prec A_2$) $\to$ DiscloseActive $\to$ Include($A_1$) & P9, P10\\
MUT14 & Idempotency Collision & CommitRequest($K_1$) $\to$ ReplayModified($K_1$) $\to$ DuplicateCommit & P1\\
MUT15 & Unauthorized Retry & RejectConflict($P$) $\to$ ImmediateRetry($P$) $\to$ UncheckedCommit & P15\\
\bottomrule
\end{tabularx}
\end{table}

\subsection{Mutant Corpus and Dual-Model Non-Equivalence Screening}
The mutation manifest was frozen before execution. To guarantee rigorous screening without bias, candidate mutants underwent blinded dual-model review using Gemini 3.6 Flash High and Claude Sonnet 4.6 against formal contract semantics. Only candidates mutually classified as executable and non-equivalent were included in the benchmark denominator ($N = 45$), completely avoiding cosmetic or unexecutable padding.

\section{Evaluation Protocol}
\subsection{Primary Endpoint}
For any strategy $T$, the mutation score is defined as:
\begin{equation}
MS(T) = \frac{K_T}{N_{\text{included}} - N_{\text{eq}}},
\end{equation}
where $K_T$ is the number of non-equivalent mutants killed by $T$, and the denominator is the frozen corpus ($N = 45$).

\subsection{Execution Protocol and Reproducibility Freeze}
The protocol frozen under run identifier \texttt{govmut-soict-2026-final-v2} executed all 45 mutants across 16 method/seed combinations on revision \texttt{ab877e04}, totaling 720 executions. Unmutated baseline preflight passed all 16 slots with zero false kills.

\section{Results}
Table~\ref{tab:res} details the comparative evaluation across the four strategies.

\begin{table}[t]
\caption{Sealed mutant-level strategy comparison ($N = 45$ non-equivalent mutants, 720 executions).}
\label{tab:res}
\centering\small
\begin{tabularx}{\textwidth}{p{0.24\textwidth}XXX}
\toprule
Strategy & Killed / 45 & Mutation Score (95\% CI) & Exact Paired Comparison vs.\ M3\\
\midrule
M0 Regression & 16 & 0.356 (0.232--0.502) & $p = 0.125$\\
M1 Stateless properties & 4 & 0.089 (0.035--0.207) & $p = 3.05 \times 10^{-5}$\\
M2 State machine & 6 & 0.133 (0.063--0.262) & $p = 1.22 \times 10^{-4}$\\
M3 Combined & 20 & 0.444 (0.309--0.588) & Reference baseline\\
\bottomrule
\end{tabularx}
\end{table}

M3 detected all 16 mutants killed by M0 and four additional mutants ($M0\text{-only} = 0, M3\text{-only} = 4$). The four additional kills achieved by M3 were complex multi-action defects (MUT8 digest tampering under race conditions, MUT12 cache re-inflation, MUT13 superseded re-selection, and MUT15 unauthorized retry escalation). While the 8.9 percentage point improvement over regression alone was not statistically significant in exact paired testing ($p = 0.125$), M3 significantly outperformed M1 ($p = 3.05 \times 10^{-5}$) and M2 ($p = 1.22 \times 10^{-4}$). Across all 720 individual method/seed executions, the terminal ledger recorded 161 \textsc{Killed}, 559 \textsc{Survived}, and zero \textsc{Infrastructure Error} events.

\section{Interpretation and Discussion}
The empirical findings demonstrate clear **complementarity** rather than unilateral replacement. Regression suites excel at guarding known single-point enforcement gates, whereas state-machine testing explores multi-hop transition orderings that expose subtle sequence-dependent vulnerabilities. Surviving mutants highlight specific areas for test suite expansion (e.g., deeper multi-actor interleaving and bitemporal conflict fuzzing) rather than benign faults.

\section{Threats to Validity}
\textbf{Construct Validity:} Fault operators directly reflect real-world healthcare governance hazards. Non-equivalence was screened via blinded dual-model evaluation; future iterations will incorporate panel adjudication.
\textbf{Internal Validity:} Fixed seeds, frozen budgets, and persistent shrunk counterexample logs ensure deterministic reproducibility.
\textbf{External Validity:} Evaluation was conducted on CLARA-Care/GLHS; replication on other agentic EHR platforms is encouraged.
\textbf{Clinical Safety Boundary:} This study evaluates software contract assurance and does not claim medical diagnosis accuracy or clinical efficacy.

\section{Reproducibility}
All evaluation artifacts, test seeds, shrunk counterexamples, and the complete mutant-by-strategy execution matrix are archived under run ID \texttt{govmut-soict-2026-final-v2} in \path{research/assurance_soict/results/final-analysis.json}.

\section{Relationship to Prior GLHS Work}
While the GLHS architecture manuscript investigates the co-versioned disclosure-to-commit protocol~\cite{glhs}, GovMut-Health addresses the orthogonal software-engineering problem of benchmarking test adequacy for sequence-sensitive governance contracts.

\section{Conclusion}
GovMut-Health provides the first domain-specific mutation benchmark for stateful healthcare AI systems. By formalizing 15 multi-action defect classes, establishing cryptographic Merkle state transition verification (Theorem~3), and comparing four testing strategies under sealed execution, the benchmark demonstrates that combining regression and generated state-machine testing achieves the highest defect coverage (44.4\%).

\begin{thebibliography}{16}
\bibitem{quickcheck} Claessen, K., Hughes, J.: QuickCheck: A lightweight tool for random testing of Haskell programs. In: ICFP 2000, pp. 268--279. ACM (2000)
\bibitem{utting} Utting, M., Pretschner, A., Legeard, B.: A taxonomy of model-based testing approaches. Software Testing, Verification and Reliability 22(5), 297--312 (2012)
\bibitem{jia} Jia, Y., Harman, M.: An analysis and survey of the development of mutation testing. IEEE Trans. Softw. Eng. 37(5), 649--678 (2011)
\bibitem{pretschner2008} Pretschner, A., Mouelhi, T., Le Traon, Y.: Model-Based Tests for Access Control Policies. In: ICST 2008, pp. 338--347. IEEE (2008)
\bibitem{xu2012} Xu, D., Thomas, L., Kent, M., et al.: A Model-Based Approach to Automated Testing of Access Control Policies. In: SACMAT 2012. ACM (2012)
\bibitem{nist800192} Hu, V.C., Kuhn, D.R., Yaga, D.: Verification and Test Methods for Access Control Policies/Models. NIST SP 800-192 (2017)
\bibitem{memtx} Li, X., Wang, Y., Lu, H., et al.: MemTX: Transactional Belief Commit for Stateful Agent Memory. arXiv:2607.23929 (2026)
\bibitem{memtxn} Cui, J., Zhang, L., et al.: MemTxn: Snapshot Isolation and Conflict Resolution for Stateful LLM Agents. arXiv:2607.24102 (2026)
\bibitem{cordon} Evans, M., et al.: Cordon: Staged Irreversible Effects in LLM Agent Transactions. arXiv:2606.18901 (2026)
\bibitem{verifiedconcurrency} Khan, S.: Verified Detection and Prevention of Concurrency Anomalies in Multi-Agent Large Language Model Systems. arXiv:2606.17182 (2026)
\bibitem{swemutation} Sun, Y., Zhao, Y., Wang, Y., et al.: SWE-Mutation: Can LLMs Generate Reliable Test Suites in Software Engineering? arXiv:2605.22175 (2026)
\bibitem{toki} Wang, Z.: TOKI: A Bitemporal Operator Algebra for Contradiction Resolution in LLM-Agent Persistent Memory. arXiv:2606.06240 (2026)
\bibitem{fhirhttp} HL7 International: FHIR R4: RESTful API / HTTP (2019), \url{https://hl7.org/fhir/R4/http.html}
\bibitem{fhirconsent} HL7 International: FHIR R4: Consent (2019), \url{https://hl7.org/fhir/R4/consent.html}
\bibitem{prov} W3C: PROV-O: The PROV Ontology. W3C Recommendation (2013)
\bibitem{glhs} Thien, N.N., Quang, T.M.: GLHS: A Co-Versioned Disclosure-to-Commit Governance Contract for Longitudinal Health AI. Working preprint (2026)
\end{thebibliography}
\end{document}
